% Section 13 of the manuscript: limitations. It follows the section plan of outline.md and sections 7 and 10.11 of the
% registration; where a limitation has an expected direction, the direction is stated as an expectation, not as a
% result.

\section{Limitations}
\label{sec:limitations}

\paragraph{The reference.} The reference is a reference input field, not the input method's internal state, and
another application may receive other calls from the same input method. The ground truth depends on the rules that
assemble sessions from the calls; with live conversion on, the rule folds the automatic commits made while a reading is
typed into one session, and a rule that did not fold them would define a conversion differently. The reference is
software on the measured machine, with no external clock or device, so the agreement is between two software records;
and the kind of field affects the reference, as the two Windows fields of the pilot showed
(Section~\ref{sec:method-reference}).
The field can lose the focus by itself, as it did once in the pilot; the presses sent during such a loss reach another
application, and their trials stay in the evaluation as disagreements between the plan and the ground truth.

\paragraph{Windows.} The Windows estimates include the difference between the two constructs
(Section~\ref{sec:observability-consequence}). They hold for an unknown mixture of the
logger's two paths, because a row does not record its path (Section~\ref{sec:logger-windows}). The input method ran
with its default settings, with predictive candidates shown; with them turned off, the window would open at the
conversion rather than while the reading is typed. Applications that draw their own candidates were not tested, and in
them the logger may see no session. The counts \texttt{nav} and \texttt{abandon} are counted system-wide while a
session is open and are upper bounds; this design does not exercise their excess, because only the executor typed. In
this study the per-text abandonment rate of the Windows logger was lower than the reference's
(Section~\ref{sec:results-session}). In
S15 the input method commits when another window comes to the front, so the session ends there. On the cloaking path
the observed interval ends at least 150~ms before the row's time, and on both paths the row's time is that of its
collection; the matching is not adjusted for either, so short sessions can fail to pair for this reason. The input mode
was switched with keys that a program sends, and the response of the input method was not compared with that to keys
pressed by hand.

\paragraph{Synthetic and idealised input.} The procedures were calibrated so that each conversion ends within its trial,
and trials were separated by a pause of 1~s; a conversion that ran from one trial into the next was excluded as
carry-over, so the agreement for conversions that continue into the input that follows is not estimated. Each key was
held for 60~ms and no two presses overlapped; the session rules of the macOS logger act on the order of the presses,
which overlapping presses do not change, and the release rows, whose times overlapping would change, were not validated.
The think time before the decisive key was drawn between 0.2 and 20~s (Section~\ref{sec:method-plans}); a difference
in where the rules and the reference end a session can include such a pause, so the long tails of the duration
differences may depend on this distribution, and the limits of agreement are not carried over to other distributions
of pauses. Key intervals were drawn independently of the keys, with no pauses at the boundaries of words and phrases;
this affects the measures of timing and the feature sets of question~4 that keep time. The timing of human typing
carries information about the keys \citep{monaco2018} that these logs lack, while the attacker's knowledge of the
generator and the absence of typing errors favour identification here, so how those sets would fare on human typing is
not known. The input had no typing errors, and
corrections occurred only where a scenario prescribed them; errors would add presses that the attacks of question~4 do
not expect, so the identification is expected to be lower with them than here, and the effect of errors on detection is
not known. The events were injected through the interfaces of the operating system, so the timing of the key rows
(Table~\ref{tab:key-timing}) covers the path from the injection to the row; the path from a physical key, which earlier
validations measured from the key to the record \citep{pinet2017}, adds its own delay, so the delay from a physical key
to a row is expected to be longer than the delays reported here. The rows take their times from the wall clock when an
event reaches the logger; the registration asks for network time synchronisation to be off during the runs, whether it
was off was not recorded (Section~\ref{sec:method-cells}), and an adjustment of the clock while the logger runs, which
in ordinary use can happen, moves the times of the rows. Such an adjustment was not tested; during the runs no trial was
excluded for a step of the wall clock (Section~\ref{sec:results-detection}), but a smaller adjustment, or one undone
within a block, need not show in the drift between the clocks of the records, which was measured only at the block
markers (Section~\ref{sec:results-reference}), so an effect of such adjustments on the times within a trial is not
ruled out. The input mode was switched only by the mode keys; on macOS a switch made
otherwise is followed by polling the input source every 100~ms, which was not exercised.

\paragraph{The rules of the macOS logger.} The minus, comma and period keys close a session, so a reading that contains
a long-vowel mark or punctuation is split into several sessions, which can raise the number of conversions that a
researcher counts from the log. With the kana layout every kana typed with a symbol key, among them the voiced-sound
mark, closes a session in the same way; S12 used only kana typed with letter, digit, minus, comma and period keys, so
its agreement does not represent kana input in general. The cut-off leaves out
the sessions that reach 20~s; how it moves each aggregate is stated in Section~\ref{sec:recommendations}.

\paragraph{One machine and one version per system.} One input method was measured per system, on one machine, with the
versions recorded in the run manifests; the Windows versions of the measurement differed from those of the pilot, and
the registered check passed (Section~\ref{sec:method-cells}). The learning of the input method was not reset, so the
candidates and the automatic commits may have changed over the runs. Other programs ran on the measuring machines
(Section~\ref{sec:results-attempts}), and their load was not controlled. Scenario S9 (choosing a candidate with the
mouse) was not evaluated.

\paragraph{The prediction.} The prediction models the macOS logger's rules and the runner's presses, but not the
modifier keys, the collection of rows or the delay from dispatch to the logger, so a difference between the prediction
and the observation is reported as such and is not attributed to a cause. On macOS as on Windows the time of a session
row is that of its collection, up to about 100~ms after the session closed, and the matching is not adjusted for it, so
short sessions can fail to pair for this reason; the strata in which the logger's rows departed most from the specified
rows (Section~\ref{sec:results-detection}) may owe their departures to this, which was not tested.

\paragraph{Deviations.} A defect of the frozen code was corrected next to the registered analysis (deviation D10), and
three changes were made after the results had been seen: two of reporting (D11 and D12) and one that describes three
feature sets of question 4 as the code computed them and adds an exploratory analysis (D13;
Section~\ref{sec:deviations}).

\paragraph{Scopes and intervals.} The primary scope sets aside the trials whose presses went into a session that is cut
off; how often that happens is a property of this design, and the trials that remain differ between speeds in reading
length and think time. The Wilson intervals of pooled proportions are read as lower bounds of the uncertainty
(Section~\ref{sec:method-analysis}), and the bootstrap intervals cover dependence within a block but not across blocks.

\paragraph{The stopping rule.} The rules that stop a run when a person's key presses may have entered its logs have
stated blind spots (Section~\ref{sec:method-checks} and section~7 of the registration): for example, a press that did
not reach the reference and a press from outside cancel out, and a person's presses that reach neither the reference nor
another application are caught by neither rule. On Windows the rule of deviation D4 can also stop a run for an error of
the logger, because a key row that no sent press explains can come from either, so the runs analysed can be selected for
the absence of such errors; had every attempt been stopped, the cell would not have been measured.

\paragraph{Recoverability.} The study judges one phrase at a time, not a whole log, and no record that keeps times;
what else it does not evaluate, and that its attacks are not claimed to be the strongest possible, is stated in
Section~\ref{sec:recoverability-method}. The Wilson interval of the near-chance rule treats trials as independent, so
with positive dependence within a block the rule can be met more easily than it should be
(Section~\ref{sec:recoverability-method}).
The decoys were not matched to the true phrase in length, so the numbers of
keys of each category, which both registered sets without times keep (deviation D13), are expected to distinguish the
candidates more than they would among phrases of equal length.

\paragraph{Runtime cost.} The values of each system are read within that system (Section~\ref{sec:cost-method}), and
battery use was not measured.
